\documentclass{article}
\usepackage{tcee-repaso}
\tcetema[Banco Mundial y bancos multilaterales de desarrollo]{3.B.30}{El Grupo del Banco Mundial, los Bancos Regionales de Desarrollo y otras instituciones financieras multilaterales de desarrollo.}
\tcefecha{19/11/2026}
\begin{document}
\cabecerarepaso
\begin{columnas}

% ==========================================================================
\bloque[2,5]{Introducción}
\begin{itemize}
\item \textbf{Enganche}: «quien a buen árbol se arrima, buena sombra le cobija».
  \begin{itemize}
  \item GBM, ejercicio 2026: \dato{123.500 M\$, el doble que hace 10~años}{BM}.
  \item Reformas de capital: \dato{+792.000 M\$ en una década}{BMD, 10/2026}.
  \item ¿Por qué un banco de Estados y cuánto puede prestar a la sombra de sus accionistas?
  \end{itemize}
\item \textbf{Relevancia}:
  \begin{itemize}
  \item Financiación larga, barata y anticíclica; la AIF para los más pobres \vertema{3.B.29}.
  \item Instrumento del G-20 para clima y ODS \vertema{3.B.19}; España accionista de todos \vertema{4.A.24}.
  \end{itemize}
\item \textbf{Contextualización}: \concepto{BMD}: institución de Estados que capta en los mercados y de los donantes y presta, con conocimiento, a países en desarrollo.
  \begin{itemize}
  \item El FMI no es de desarrollo \vertema{3.B.22}; Bretton Woods \vertema{3.B.20}; regionales 1959-1966; BERD 1991; BAII y NDB, 2015-2016.
  \end{itemize}
\item \textbf{Problemática}: ¿por qué existen y cuánto pueden prestar? ¿Cómo es el GBM? ¿Qué aportan los demás y cómo se reforma el sistema?
\item \textbf{Estructura}: lo común (\hyperref[rep:I]{I}), el mayor (\hyperref[rep:II]{II}), regionales, nuevos y reforma (\hyperref[rep:III]{III}).
\end{itemize}
\esquemacante{%
\apartadocante{Introducción}{2,5}
\apartadocante{I. Los BMD: elementos comunes}{8,5}
\subcante{1. Definición y razón de ser}{5}
\subcante{2-5. Rasgos, tipos, gobierno, actuación}{3,5}
\apartadocante{II. El Grupo del Banco Mundial}{9,5}
\subcante{1. Información general e historia}{2,5}
\subcante{2. Estructura y estrategia}{2}
\subcante{3. BIRF, AIF, CFI, MIGA, CIADI}{5}
\apartadocante{III. Regionales y otras}{8}
\subcante{1. Bancos regionales}{4,5}
\subcante{2. Subregionales y temáticas}{0,5}
\subcante{3. Nuevas: BAII y NDB}{1}
\subcante{4. Reforma del sistema y España}{2}
\apartadocante{Conclusión}{1,5}}

% ==========================================================================
\bloque[8,5]{I. Los BMD: elementos comunes\label{rep:I}}
\begin{ideaclave}
Un BMD presta barato y a largo plazo a quien el mercado no financia gracias a la garantía de sus accionistas; cuánto presta lo decide su capital, y para qué, sus votos.
\end{ideaclave}
\sub{I.1. Definición y razón de ser\label{rep:I.1}}
\subsub{a) Intermediación con la garantía de los accionistas}
\begin{itemize}
\item \emph{Idea}: fallos del crédito soberano: sin tribunal para Estados, información asimétrica, pecado original \vertema{3.B.29}.
\item \emph{Supuestos}:
  \begin{itemize}
  \item patrimonio $E$ (desembolsado + reservas); préstamos $L=E+B$;
  \item \concepto{capital exigible} (suscrito, no desembolsado): garantía, no recurso para prestar;
  \item se endeuda a $r_B$ (AAA), presta a $r_B+m$; mínimo $E/L\geq e^{\min}$.
  \end{itemize}
\item \emph{Resultado}, la palanca:
\funcionclave{L^{\max}=\frac{E}{e^{\min}}}
  \begin{itemize}
  \item Ahorro del prestatario: $(r^*+\rho)-(r_B+m)$.
  \item Beneficio $r_BE+(m-c)L$; sin capital nuevo, $g_L=r_B+(m-c)/e^{\min}$.
  \end{itemize}
\item \emph{Implicaciones}:
  \begin{itemize}
  \item 1 \$ de patrimonio $\to$ $1/e^{\min}$ de préstamos (20\,\%: 5; 18\,\%: 5,6);
  \item 3~palancas: $E$, $e^{\min}$, $m$; dilema préstamos frente a AAA;
  \item \concepto{acreedor preferente}: le pagan antes que a otros.
  \end{itemize}
\end{itemize}
\graficorepaso[0.9\linewidth]{palanca-capital}{Bajar $e^{\min}$ mueve a lo largo de la curva; el capital híbrido la desplaza}
\begin{itemize}
\item \emph{Extensiones}: revisión de los marcos de capital del G-20 (2022): apetito de riesgo, capital exigible, transferir riesgo, agencias.
  \begin{itemize}
  \item \concepto{Capital híbrido}: perpetuo y subordinado; cuenta como patrimonio.
  \item BIRF: 20 $\to$ 19\,\% (2023) $\to$ 18\,\% (2024); híbrido de 12~accionistas \dato{1.200 M\$}{BM}; BAfD, primer híbrido en el mercado (2024); FMI: comprarlo con DEG.
  \end{itemize}
\item \emph{Evidencia}: impagos muy bajos (\autor{Perraudin} et al., 2016); las agencias apenas valoran el capital exigible (\autor{Humphrey}, 2017).
\item \emph{Valoración}: riesgo de cambio en el país; no explica los bienes globales.
\end{itemize}
\subsub{b) Bienes públicos globales (\autor{Kaul} et al., 1999)}
\begin{itemize}
\item \emph{Idea}: no rivales ni excluibles entre países: clima, pandemias, estabilidad financiera.
\item \emph{Modelo}: $N$ países, $G=\sum g_i$, bienestar $b_i(G)-c\,g_i$.
  \begin{itemize}
  \item Nash: $b_i'(G^N)=c$ (polizón) \vertema{3.A.23}; óptimo (\autor{Samuelson}, 1954):
  \end{itemize}
\funcionclave{\sum_{i=1}^N b_i'(G^*)=c}
  \begin{itemize}
  \item $G^*>G^N$; la brecha crece con $N$.
  \end{itemize}
\end{itemize}
\graficorepaso[0.85\linewidth]{bienes-publicos-globales}{Suma vertical de beneficios marginales: $G^*$ frente a $G^N$}
\begin{itemize}
\item \emph{Implicación}: donaciones y fondos, no préstamos.
\item \emph{Extensiones}: agregación (\autor{Hirshleifer}, 1983): suma, eslabón más débil, mejor disparo.
\item \emph{Evidencia}: clima de los BMD a renta baja y media \dato{102.600 M\$}{2025}.
\item \emph{Valoración}: tensión con la pobreza («deriva del mandato», EE.\,UU., 2025).
\end{itemize}
\subsub{c) Efecto catalizador (\autor{Rodrik}, 1995)}
\begin{itemize}
\item \emph{Modelo}: buenos ($p_A$) y malos ($p_B$), proporción $\lambda$; $\bar p=\lambda p_A+(1-\lambda)p_B$; el mercado se cierra si
\funcionclave{p_A R\geq1+r^*>\bar p R}
\item El banco examina, cofinancia y revela el tipo (\autor{Akerlof}; \autor{Stiglitz}-\autor{Weiss}) \vertema{3.A.13}.
\item \emph{Extensiones}: \concepto{ratio de movilización} (privado por dólar del BMD); \concepto{adicionalidad} (financiar lo que el mercado no); condicionalidad; versión macro: FMI \vertema{3.B.22}.
\item \emph{Evidencia}: «de miles de millones a billones» (2015) no se cumplió: \dato{108.700 M\$ movilizados}{2024}; 0,6 \$ por dólar.
\item \emph{Valoración}: \autor{Wapenhans} (1992): cultura de aprobar.
\end{itemize}
\sub{I.2. Rasgos generales\label{rep:I.2}}
\begin{itemize}
\item Accionistas Estados, sin lucro; capital suscrito = desembolsado + exigible.
\item Fondos: bonos (primer \emph{swap}, BM-IBM, 1981); \concepto{reposiciones} de donantes; fondos fiduciarios.
\item Préstamos de mercado o concesionales; casamiento de plazos; nada a quien está en atrasos.
\end{itemize}
\sub{I.3. Tipología: las ventanillas\label{rep:I.3}}
{\small\noindent\begin{tabular}{@{}>{\raggedright\arraybackslash}p{0.17\linewidth}|>{\raggedright\arraybackslash}p{0.34\linewidth}|>{\raggedright\arraybackslash}p{0.2\linewidth}|>{\raggedright\arraybackslash}p{0.2\linewidth}@{}}
\textbf{Ventanilla} & \textbf{A quién} & \textbf{Fondos} & \textbf{Ejemplos}\\ \hline
Dura & Estados de renta media & bonos y capital & BIRF, BID, BAsD\\
Blanda & Estados pobres: créditos de favor y donaciones & reposiciones (y bonos en la AIF) & AIF, FAD, FAsD\\
Privada & empresas, sin garantía estatal & bonos y capital & CFI, BID Invest, BERD\\
Garantías & riesgo político & capital y primas & MIGA\\
\end{tabular}\par}
\begin{itemize}
\item \concepto{Graduación}: de la blanda a la dura (inversa, al revés); países \emph{blend}.
\end{itemize}
\sub{I.4. Gobierno\label{rep:I.4}}
\begin{itemize}
\item 3~niveles, casi iguales en todos:
  \begin{enumerate}
  \item \concepto{Junta de Gobernadores}: un gobernador por país (ministro de Economía); 1~vez al año; capital, miembros, convenio;
  \item \concepto{Directorio Ejecutivo}: gestión diaria, aprueba préstamos y estrategias; directores por país o grupo (\emph{silla});
  \item \concepto{presidente}: dirige al personal.
  \end{enumerate}
\item Votos: igualdad entre Estados + peso del capital:
\begin{formula}
\[\text{votos}_i=\text{votos básicos}+\text{acciones suscritas}_i\]
\end{formula}
\item Estrategias corporativa, sectoriales y de país.
\item Equilibrio: en los regionales, la región conserva la mayoría; en el BM, más votos a los países en desarrollo (2010).
\end{itemize}
\sub{I.5. Actuaciones\label{rep:I.5}}
\begin{itemize}
\item \emph{Financiación}:
  \begin{itemize}
  \item proyectos de inversión, sobre todo infraestructuras;
  \item apoyo presupuestario ligado a reformas (se desembolsa al cumplir);
  \item programas por resultados (pagan al medir resultados, no al gastar).
  \end{itemize}
\item \emph{Ciclo del proyecto}: (1) identificación y evaluación; (2) licitación abierta a los miembros; (3) ejecución y supervisión; (4) evaluación independiente.
\item \emph{Salvaguardias} (reasentamientos, pueblos indígenas, biodiversidad): encarecen, pero evitan daños; Marco Ambiental y Social (2018).
\item \emph{Asistencia técnica} y \emph{conocimiento} (\emph{Informe sobre el desarrollo mundial}, estadísticas de deuda).
\item \emph{Resultados}: marcos de resultados desde París (2005) y Busan (2011); evaluación que responde al Directorio; Panel de Inspección (1993); evaluaciones aleatorias \vertema{3.B.28}.
\item Un BMD es a la vez banco, fondo de donantes y agencia de conocimiento.
\end{itemize}

% ==========================================================================
\bloque[9,5]{II. El Grupo del Banco Mundial\label{rep:II}}
\begin{ideaclave}
De reconstruir Europa a la pobreza y, desde 2023, a un «planeta habitable»; 5~entidades, una por ventanilla; la reforma reciente: prestar más con el mismo capital.
\end{ideaclave}
\sub{II.1. Información general e historia\label{rep:II.1}}
\begin{itemize}
\item BIRF 1944 (189~miembros), CFI 1956 (186), AIF 1960 (175), CIADI 1966 (159), MIGA 1988 (182). Antes, ser miembro del FMI.
\item \emph{Etapas}:
  \begin{itemize}
  \item 1947: Francia, 250~M\$; AAA en 1959; 60: AIF;
  \item 70: \autor{McNamara}, pobreza; informe \autor{Pearson} (1969);
  \item 80: ajuste estructural, Consenso de Washington \vertema{3.B.28};
  \item 90: \autor{Wapenhans}, Panel de Inspección, \autor{Wolfensohn}, PPME;
  \item 2000: ODM, IADM; 2010: voz y capital, objetivos gemelos (2013), capital (2018);
  \item 2020: hoja de ruta de la evolución, garantías en MIGA, 22~indicadores, Misión 300.
  \end{itemize}
\end{itemize}
\graficorepaso[0.9\linewidth]{compromisos-gbm}{Compromisos del GBM, ejercicios 2016-2026: de 61.000~a 123.500~M\$}
\sub{II.2. Estructura y estrategia\label{rep:II.2}}
\begin{itemize}
\item Junta, Directorio de 25, presidente estadounidense (\autor{Ajay Banga}, desde 2023), Comité para el Desarrollo con el FMI.
\item \emph{Votos}: 5,55\,\% básicos.
  \begin{itemize}
  \item EE.\,UU. \dato{15,98\,\%}{veto con el 85\,\%}; España \dato{1,84\,\%}{puesto 14}, silla con México, Venezuela y Centroamérica (4,59\,\%).
  \item 2010: +86.200 M\$ y 47,19\,\% para los países en desarrollo; 2018: 13.000~M desembolsados; 2025: sin aumento selectivo (hace falta el 75\,\%).
  \end{itemize}
\end{itemize}
\graficorepaso[0.95\linewidth]{gobierno-gbm}{El gobierno del GBM y el voto en el BIRF}
\graficorepaso[0.85\linewidth]{votos-birf}{Votos en el BIRF: EE.\,UU., único con veto}
\begin{itemize}
\item \emph{Estrategia}:
  \begin{itemize}
  \item objetivos gemelos (2013): pobreza extrema (3 \$/día desde 2025) < 3\,\% y renta del 40\,\% más pobre;
  \item misión 2023: «en un planeta habitable»;
  \item hoja de ruta: patrimonio mínimo 20 $\to$ 18\,\%, capital híbrido y garantías; 22~indicadores; 6~programas de desafíos globales (2024);
  \item empleo, prioridad desde 2025 (5~sectores, \dato{55\,\%}{2026}); clima \dato{44\,\%}{2026} de la financiación propia.
  \end{itemize}
\end{itemize}
\sub{II.3. Las 5~entidades\label{rep:II.3}}
\begin{itemize}
\item \emph{Clasificación por renta} (ej.\ 2027, Atlas): baja $\leq$1.175 \$; media-baja $\leq$4.635; media-alta $\leq$14.375; alta.
\end{itemize}
\subsub{a) BIRF}
\begin{itemize}
\item Estados de renta media; \dato{332.500 M\$ suscritos, 23.600~desembolsados}{BM}; bonos AAA.
\item SOFR + margen, hasta 35~años (medio 20); \dato{50.200 M\$}{2026, récord}; $E/L$ 21,6\,\% (mínimo 18).
\item Críticas: EE.\,UU., renta media, lentitud; préstamos en crisis (terreno del FMI).
\end{itemize}
\subsub{b) AIF}
\begin{itemize}
\item Umbral \dato{1.365 \$}{ej.\ 2027}; 78~países (19~\emph{blend}); asignación por desempeño.
\item Créditos a 31~años al 0,75\,\% o 40~al 0\,\%; donaciones si hay riesgo de deuda; bonos desde 2018.
\item \textbf{AIF21}: \dato{100.000 M\$ con 23.700~de donantes}{Seúl, 12/2024}; España \dato{400 M€ (+37\,\%)}{España}.
\end{itemize}
\subsub{c) CFI}
\begin{itemize}
\item Sector privado, sin garantía estatal; hasta 25\,\% del coste, nunca mayoritaria; préstamos B.
\item \dato{105.600 M\$ (32.000~propios)}{2026}.
\end{itemize}
\subsub{d) MIGA}
\begin{itemize}
\item Seguro contra riesgo político; \dato{10.600 M\$}{2026}; plataforma de garantías del grupo.
\end{itemize}
\subsub{e) CIADI}
\begin{itemize}
\item Conciliación y arbitraje (Convenio de Washington, 1965); \dato{63 casos en 2025}{CIADI}.
\item Críticas y salidas: Bolivia, Venezuela; Ecuador volvió (2021), Honduras (2026).
\item España: \dato{44 casos}{5.º demandado}, muchos por las renovables.
\end{itemize}

% ==========================================================================
\bloque[8]{III. Regionales y otras instituciones\label{rep:III}}
\begin{ideaclave}
Los regionales copian al BIRF más cerca de los prestatarios; los nuevos reflejan a los emergentes; el sistema aprende a prestar más con el mismo capital.
\end{ideaclave}
\sub{III.1. Los bancos regionales\label{rep:III.1}}
\begin{itemize}
\item \emph{Origen}: motivos económicos (recursos y productos a medida), políticos (voz tras la descolonización) y estratégicos (EE.\,UU., Japón).
\item Prestan solo a la región, que tiene la mayoría del capital.
\end{itemize}
{\footnotesize\noindent\begin{tabular}{@{}>{\raggedright\arraybackslash}p{0.13\linewidth}|>{\raggedright\arraybackslash}p{0.1\linewidth}|>{\raggedright\arraybackslash}p{0.37\linewidth}|>{\raggedright\arraybackslash}p{0.17\linewidth}|>{\raggedright\arraybackslash}p{0.1\linewidth}@{}}
\textbf{Banco} & \textbf{Año} & \textbf{Rasgo} & \textbf{2025} & \textbf{España}\\ \hline
BID & 1959 & 48 (26 prestatarios); BID Invest, BID Lab; Madrid 1974 & 17.881 M\$ & 1,96\\
BAfD & 1964 & 81 (54); FAD-17: 11.000 M\$; híbrido 2024 & 11.000 M\$ & 1,02 c\\
BAsD & 1966 & 69 (50); fondo blando solo donaciones; sin límite de préstamos (2026) & 29.300 M\$ & 0,56\\
BERD & 1991 & mandato político; $\geq$60\,\% privado; +4.000 M€ (2023) & 16.800 M€ & 3,65 c\\
BEI & 1958 & banco de la UE; BEI Global & 100.000 M€ & 11,27 c\\
\end{tabular}\par}
\begin{itemize}
\item c: \% del capital. Presidencias: \autor{Goldfajn} (BID), \autor{Ould Tah} (BAfD), \autor{Kanda} (BAsD), \autor{Renaud-Basso} (BERD), \autor{Calviño} (BEI).
\end{itemize}
\sub{III.2. Subregionales y temáticas\label{rep:III.2}}
\begin{itemize}
\item \emph{Subregionales y de afinidad}: CAF (1970; récord \dato{$\approx$18.500~M\$}{2025}; España, accionista de la serie C), BCIE, BDCE, BIsD.
\item \emph{Fondos verticales}: FMAM, Fondo Verde (\dato{3.260 M\$}{2025}) \vertema{3.B.31}, FIDA. ONU sin préstamos: UNCTAD, PNUD.
\end{itemize}
\sub{III.3. Nuevas: BAII y NDB\label{rep:III.3}}
\begin{itemize}
\item \textbf{BAII} (2016, Pekín): 111~miembros; China 26,4\,\% y veto (75\,\%); España fundadora, 1,72\,\%; \autor{Zou Jiayi} desde 2026; 10.590~M\$.
\item \textbf{NDB} (BRICS, 2015, Shanghái): 10~miembros; fundadores a partes iguales; moneda local; \autor{Rousseff}; 3.171~M\$.
\item Frente al BIRF: impulso de China y BRICS; foco en infraestructuras; el BAII copia normas y cofinancia con el BM.
\item Bancos chinos bilaterales \vertema{3.B.29}.
\end{itemize}
\sub{III.4. Reforma del sistema y España\label{rep:III.4}}
\subsub{a) La reforma}
\begin{itemize}
\item \emph{Mandato}: Agenda 2030 y Addis Abeba (2015).
\item \emph{Triple agenda} (\autor{Summers}-\autor{Singh}, 2023): triple mandato; triplicar hasta \dato{$\approx$390.000~M\$/año}{2030}; movilizar privado.
  \begin{itemize}
  \item Hoja de ruta del G-20 (Río, 2024); \dato{+792.000 M\$ en una década}{10/2026}; Sevilla (2025).
  \item EE.\,UU.: volver a la misión, menos clima, graduar a China.
  \end{itemize}
\end{itemize}
\subsub{b) El papel de España}
\begin{itemize}
\item Gobernador: el ministro de Economía; Tesoro (SG de Bancos Multilaterales de Desarrollo, FMI y G-20).
\item SEC, ICEX y oficinas comerciales para las licitaciones \vertema{4.A.24}.
\item FEDES (Ley 1/2023, AECID) \vertema{4.A.26}; BERD 137~M€ y BID Invest 140~M\$.
\end{itemize}
\graficorepaso[0.9\linewidth]{votos-espana}{España en los BMD (\% de votos o de capital)}

% ==========================================================================
\bloque[1,5]{Conclusión}
\begin{itemize}
\item \textbf{Recapitulación}:
  \begin{itemize}
  \item 3~razones de ser: garantía de los accionistas, bienes globales, catalizador;
  \item $L^{\max}=E/e^{\min}$; el GBM como modelo; regionales y nuevos.
  \end{itemize}
\item \textbf{Relevancia}: la palanca más barata: más de 5 \$ de préstamos por dólar de capital.
\item \textbf{Extensiones}: \vertema{3.B.22} \vertema{3.B.29} \vertema{3.B.28} \vertema{3.B.31} \vertema{4.A.24} \vertema{4.A.26}.
\item \textbf{Opinión}: reforma financiera hecha; faltan la voz y la operativa (rapidez, moneda local).
\item \textbf{Idea final}: el país se arrima al árbol de los accionistas; pero la sombra depende del árbol: ¿lo dejarán crecer, y decidir dónde da a quienes viven bajo él?
\end{itemize}

% ==========================================================================
\anexorepaso
\bloque{Datos clave}
\begin{itemize}
\item GBM: \dato{123.500 M\$}{ej.\ 2026}; BIRF \dato{50.200}{récord}; movilización \dato{112.000 M\$}{BM}.
\item BIRF: $E/L$ \dato{21,6\,\%}{mínimo 18}.
\item AIF21: \dato{100.000 M\$}{23.700 de donantes}.
\item Clima BMD: \dato{102.600 M\$}{2025}.
\item Reformas de capital: \dato{+792.000 M\$ en 10~años}{BMD, 2026}.
\end{itemize}
\bloque{Autores y fechas}
\hito{1944}{Bretton Woods: FMI y BIRF}
\hito{1954}{\autor{Samuelson}: bienes públicos}
\hito{1960}{AIF}
\hito{1983}{\autor{Hirshleifer}: eslabón más débil}
\hito{1992}{Informe \autor{Wapenhans}}
\hito{1995}{\autor{Rodrik}: catalizador}
\hito{1999}{\autor{Kaul} et al.: BPG}
\hito{2016}{BAII}
\hito{2023}{\autor{Banga}; triple agenda}
\hito{2024}{BIRF 18\,\%; AIF21}
\bloque{Preguntas del tribunal}
\begin{enumerate}
\item \textbf{¿Por qué un BMD si hay mercados?} Crédito soberano, bienes globales, información.
\item \textbf{¿Qué limita lo que presta?} $E$ y el $e^{\min}$ de su AAA.
\item \textbf{¿Por qué el capital exigible no basta?} Nunca se ha exigido.
\item \textbf{¿AIF o BIRF?} Renta bajo el umbral y sin solvencia: AIF; una sola, \emph{blend}.
\item \textbf{¿Funcionó «de miles de millones a billones»?} No: 0,6 \$ por dólar.
\end{enumerate}
\begin{trampa}
\begin{itemize}
\item AIF = Asociación Internacional de \emph{Fomento}.
\item Mínimo del BIRF: 18\,\% (no 20).
\item BAfD opera desde 1964; BEI desde 1958.
\item España en la CAF: serie C.
\end{itemize}
\end{trampa}

\end{columnas}
\end{document}
